% Final group report (đề bài cuối kỳ, nhóm 10). Layout shared with the other
% reports of the course in ../../shared/.
% Compile with XeLaTeX: latexmk -xelatex bao-cao-cuoi-ky-nhom10.tex
\documentclass[12pt,a4paper]{report}

\newcommand{\sharedir}{../../shared}
\newcommand{\figdir}{../outputs/figures}
\input{\sharedir/layout.tex}

% ============ TITLE PAGE FIELDS ============
\newcommand{\reporttitle}{BÁO CÁO CUỐI KỲ}
\newcommand{\reportname}{Hồi quy tuyến tính Bayes \\
  cho hiệu suất nhiên liệu của xe mới tại Mỹ năm 1993}
\newcommand{\studentlabel}{Nhóm 10 thực hiện}
\newcommand{\studentname}{Nguyễn Thục Anh & 25007949\\
  Bùi Quang Huy & 25007877\\
  Nguyễn Thị Bích Ngọc & 25007937\\
  Nguyễn Thu Ngân & 25007865\\
  Nguyễn Hồng Phúc & 25007860}

\begin{document}

\makereporttitlepage

\chapter*{Tóm tắt}
\addcontentsline{toc}{chapter}{Tóm tắt}

Nhóm dùng hồi quy tuyến tính Bayes để giải thích và dự đoán hiệu suất nhiên liệu
trong thành phố (biến City, đơn vị MPG) của 93 mẫu xe mới bán tại Mỹ năm 1993,
dựa trên công suất, trọng lượng và xuất xứ. Nhóm xây dựng năm mô hình, ước lượng
bằng NUTS trong PyMC, kiểm tra hội tụ, so sánh bằng PSIS-LOO và WAIC, và đối
chiếu với công thức hậu nghiệm Normal--Inverse-Gamma cho mô hình đơn giản nhất.
Kết quả cho thấy trọng lượng là biến liên hệ mạnh nhất với City: khi thêm trọng
lượng vào mô hình, hệ số của công suất gần như bằng 0, còn mỗi \SI{100}{kg} nặng
hơn ứng với City thấp hơn 6{,}2\%. Chênh lệch City giữa xe Mỹ và xe nước ngoài
cũng giảm nhiều khi có trọng lượng. Đây là liên hệ trên dữ liệu quan sát, không
phải quan hệ nhân quả. Mô hình cuối dùng hàm hợp lý Student-t để giảm ảnh hưởng
của vài xe cỡ nhỏ tiết kiệm bất thường. Với một xe Mỹ 300 hp nặng
\SI{1681}{kg}, mô hình dự đoán trung vị City khoảng 16{,}6 MPG, khoảng dự đoán
95\% từ 13{,}4 đến 20{,}6 MPG; ba xe thật gần 300 hp (có trong dữ liệu học) đều
nằm trong khoảng này. Báo cáo gồm sáu chương, trong đó bốn chương chính ứng với
bốn yêu cầu (i)--(iv) của đề.

\tableofcontents
\pagebreak

\listoffigures
\listoftables
\newpage

% ===========================================================================
\chapter{Giới thiệu}
\label{sec:intro}

\section{Bối cảnh}

Thập niên 1990, ngành ô tô Mỹ làm ăn có lãi. Thu nhập tăng nên người mua chuộng
xe lớn và mạnh hơn: xe thể thao, sedan cỡ lớn, xe van và SUV. Xe mạnh và nặng
thường tốn nhiên liệu hơn, nên câu hỏi công suất và trọng lượng ảnh hưởng thế
nào đến mức tiêu hao nhiên liệu có ý nghĩa với cả người mua lẫn nhà sản xuất. Bộ
dữ liệu của đề cho phép trả lời câu hỏi này trên một mẫu ngẫu nhiên các xe bán
ra năm 1993.

\section{Mục tiêu và câu hỏi nghiên cứu}

Biến phản hồi chính là City, biến giải thích chính là Horsepower; Weight và USA
được thêm dần qua năm mô hình. Bốn chương chính ứng với bốn yêu cầu (i)--(iv) của
đề và trả lời bốn câu hỏi trong tài liệu đính kèm, như ở
bảng~\ref{tab:questions}.

\begin{table}[tp]
  \centering
  \caption{Bốn câu hỏi của đề và cách trả lời bằng hồi quy Bayes.}
  \label{tab:questions}
  \begin{tabular}{l>{\raggedright}p{5.2cm}>{\raggedright}p{7.6cm}l}
    \toprule
    \# & Câu hỏi của đề & Cách trả lời & Mục \\
    \midrule
    Q1 & Quan hệ giữa các biến, so sánh các biểu đồ phân tán
       & Ma trận biểu đồ phân tán, hệ số hậu nghiệm và HDI 95\% & \ref{sec:eda}, \ref{sec:q1} \\
    Q2 & Khác biệt giữa xe Mỹ và xe nước ngoài
       & Biến giả USA, tương tác $\HP \times \USA$, xác suất $P(\beta_{\USA} > 0 \mid y)$
       & \ref{sec:candidates}, \ref{sec:q2} \\
    Q3 & Các điểm ngoại lai và lời giải thích
       & Kiểm tra dự đoán hậu nghiệm, phần dư chuẩn hoá, Pareto $k$ của PSIS-LOO, hàm hợp lý Student-t
       & \ref{sec:q3} \\
    Q4 & Chọn biến giải thích và phản hồi; dự đoán xe 300 hp
       & Phân phối dự đoán hậu nghiệm tại $\HP = 300$, so với ba xe thật gần 300 hp
       & \ref{sec:q4} \\
    \bottomrule
  \end{tabular}
\end{table}

\section{Vì sao chọn cách tiếp cận Bayes}

Với cách tiếp cận Bayes, mỗi hệ số có cả một phân phối hậu nghiệm chứ không chỉ
một giá trị ước lượng. Nhờ vậy câu hỏi như ``xe Mỹ có kém tiết kiệm hơn không''
được trả lời trực tiếp bằng một xác suất. Tiên nghiệm ít thông tin loại bỏ các
giá trị vô lý, ví dụ City âm hoặc trên 80 MPG, nhưng vẫn để dữ liệu quyết định
kết quả. Phân phối dự đoán hậu nghiệm tính đến cả việc chưa biết chắc tham số
lẫn sai số của từng quan sát. Điều này đặc biệt cần cho câu hỏi Q4, vì điểm cần
dự đoán nằm ở rìa dữ liệu. Phần lý thuyết dựa theo
\cite{gelman2013bda, mcelreath2020rethinking}.

% ===========================================================================
\chapter{Dữ liệu sử dụng}
\label{sec:data}

\section{Nguồn và cách lấy mẫu}

Bộ dữ liệu gồm 93 mẫu xe mới bán tại Mỹ năm 1993, lấy ngẫu nhiên từ
\emph{The 1993 Cars -- Annual Auto Issue} của Consumer Reports
\cite{consumerreports1993} và \emph{PACE New Car \& Truck 1993 Buying Guide}
\cite{pace1993}. Mỗi dòng là một mẫu xe.

\section{Danh sách biến}

File có 25 cột gồm tên xe, 20 biến liên tục và 4 biến mã hoá, như ở
bảng~\ref{tab:variables}. Bốn biến dùng trong mô hình được in đậm.

\begin{table}[tp]
  \centering
  \caption{Các nhóm biến trong bộ dữ liệu gốc.}
  \label{tab:variables}
  \begin{tabular}{l>{\raggedright\arraybackslash}p{12.5cm}}
    \toprule
    Nhóm & Biến \\
    \midrule
    Giá (nghìn USD) & Min Price, Mid Price, Max Price \\
    Nhiên liệu (MPG) & \textbf{City}, Open Road \\
    Động cơ & \textbf{Horsepower} (hp), Engine Size (L), Cylinders, RPM tại công suất cực đại,
      RPM khi chạy đều, Fuel Tank (L) \\
    Kích thước & \textbf{Weight} (kg), Length, Wheelbase, Width (cm), U-turn (m),
      Rear seat (cm), Luggage (m$^3$), Passengers \\
    Mã hoá & Type 1--6 (Small, Midsize, Large, Compact, Sporty, Van); Airbag 0--2;
      Drive train 0--2; Manual 0/1; \textbf{USA} 0/1 (1 là xe Mỹ) \\
    \bottomrule
  \end{tabular}
\end{table}

\section{Làm sạch dữ liệu}
\label{sec:cleaning}

Dấu \texttt{*} trong file gốc được đổi thành giá trị thiếu. Luggage capacity
thiếu 11 dòng (9 xe van cùng Chevrolet Corvette và Mazda RX-7), rear seat room
thiếu 2 dòng (Corvette, RX-7), Cylinders thiếu 1 dòng (RX-7, xe này dùng động cơ
rotary nên không có xi-lanh theo nghĩa thông thường). Hai cột RPM được lưu dạng
chuỗi có dấu phẩy ngăn cách hàng nghìn, ví dụ ``6,300'', nên được chuyển sang
số. Bốn biến của mô hình là City, Horsepower, Weight và USA không thiếu giá trị
nào, nên giữ lại cả 93 dòng; kết quả lưu trong \texttt{data\_clean.csv}.

\section{Đơn vị nhiên liệu}
\label{sec:units}

Tiêu đề cột ghi ``km/100L'' còn tài liệu đi kèm ghi ``litres per 100km'', hai
đơn vị này ngược nhau. Giá trị chạy từ 15 đến 46, và xe nhỏ nhẹ nhất là Geo
Metro (55 hp, \SI{769}{kg}) lại có giá trị cao nhất, nên cả hai cách ghi đều
không hợp lý. Để xác định đơn vị, nhóm so từng dòng với bộ dữ liệu gốc
\texttt{Cars93} trong gói R \texttt{MASS} \cite{lock1993cars, mass2002cars93}, bộ
này có cùng 93 mẫu xe theo cùng thứ tự. Kết quả:

\begin{itemize}
  \item City trùng khớp cả 93 dòng với \texttt{MPG.city}, Open Road trùng khớp
    với \texttt{MPG.highway}; tài liệu của \texttt{MASS} ghi rõ hai biến này là
    dặm trên gallon Mỹ theo xếp hạng EPA;
  \item Horsepower trùng khớp cả 93 dòng, còn Weight bằng \texttt{Weight} (pound)
    nhân $0{,}4536$, sai lệch không quá \SI{0.5}{kg} do làm tròn. Ví dụ Acura
    Integra: City 25, Open Road 31, 140 hp, \SI{1227}{kg} so với 25, 31, 140 hp
    và 2705 lb.
\end{itemize}

Vậy bản Q\_10 là \texttt{Cars93} đã đổi một số cột sang hệ mét, nhưng giữ nguyên
hai cột nhiên liệu và ghi nhầm đơn vị. Từ đây báo cáo dùng MPG (dặm trên gallon,
càng cao càng tiết kiệm). Nếu cần đổi đơn vị thì
$\text{L/100km} = 235{,}2 / \text{MPG}$.

\section{Thống kê mô tả}

Mẫu có 48 xe Mỹ và 45 xe nước ngoài. Công suất từ 55 đến 300 hp, trung vị 140;
City từ 15 đến 46 MPG, trung vị 21 (bảng~\ref{tab:describe}). Hệ số tương quan
giữa HP và City là $-0{,}67$, giữa City và Open Road là $0{,}94$.

\begin{table}[tp]
  \centering
  \caption{Thống kê mô tả các biến chính, $n = 93$.}
  \label{tab:describe}
  \begin{tabular}{lrrrrr}
    \toprule
    Biến & Trung bình & Độ lệch chuẩn & Nhỏ nhất & Trung vị & Lớn nhất \\
    \midrule
    City (MPG)            & \num{22.4} & \num{5.6}  & 15  & 21   & 46 \\
    Open Road (MPG)       & \num{29.1} & \num{5.3}  & 20  & 28   & 50 \\
    Horsepower (hp)       & \num{143.8} & \num{52.4} & 55 & 140  & 300 \\
    Weight (kg)           & \num{1394} & \num{268}  & 769 & \num{1379} & \num{1862} \\
    Engine Size (L)       & \num{2.7}  & \num{1.0}  & \num{1.0} & \num{2.4} & \num{5.7} \\
    Mid Price (nghìn USD) & \num{19.5} & \num{9.7}  & \num{7.4} & \num{17.7} & \num{61.9} \\
    \bottomrule
  \end{tabular}
\end{table}

Chia theo xuất xứ, xe Mỹ nặng hơn (trung bình \SI{1449}{kg} so với
\SI{1335}{kg}), máy lớn hơn (\SI{3.1}{L} so với \SI{2.2}{L}) và kém tiết kiệm hơn
(City \num{21.0} so với \num{23.9} MPG), trong khi công suất trung bình gần bằng
nhau (\num{147.5} so với \num{139.9} hp). Vì vậy chênh lệch City giữa hai nhóm có
thể do trọng lượng chứ chưa chắc do công suất; mục~\ref{sec:q2} sẽ kiểm tra điều
này.

\section{Phân tích khám phá}
\label{sec:eda}

Ma trận biểu đồ phân tán ở hình~\ref{fig:scatter} cho thấy City giảm theo HP,
Weight và Engine Size theo dạng cong: giảm nhanh ở nhóm xe nhỏ và chậm lại ở nhóm
xe lớn. Xe nước ngoài tập trung ở vùng nhẹ, máy nhỏ, City cao; xe Mỹ nặng hơn và
máy lớn hơn. City và Open Road gần như nằm trên một đường thẳng, nên hai biến
nhiều khả năng có xu hướng giống nhau; báo cáo chỉ xây dựng mô hình cho City.

\resultfig[0.95]{fig_2_6_scatter}{Ma trận biểu đồ phân tán các biến chính, tô màu
  theo xuất xứ.}{fig:scatter}

Theo loại xe (hình~\ref{fig:box-type}), nhóm Small tiết kiệm nhất, còn Van và
Large tốn nhiên liệu nhất. Một hồi quy OLS sơ bộ $\log(\City) \sim \HP$ cho thấy
năm xe có phần dư studentised vượt $\pm 2$: Honda Civic ($+3{,}6$), Geo Metro
($+3{,}2$), Suzuki Swift ($+2{,}4$), Ford Aerostar ($-2{,}4$) và Volkswagen
Eurovan ($-2{,}3$). Mục~\ref{sec:q3} xét lại các xe này với mô hình đầy đủ.

\resultfig[0.75]{fig_2_6_box_type}{City theo loại xe; điểm tô màu theo xuất
  xứ.}{fig:box-type}

% ===========================================================================
\chapter{Quá trình mô hình hoá}
\label{sec:modeling}

\section{Cơ sở lý thuyết hồi quy tuyến tính}

Mô hình hồi quy tuyến tính biểu diễn biến phản hồi qua các biến giải thích và một
sai số ngẫu nhiên:
%
\begin{equation}
  \label{eq:linreg}
  y_i = \beta_0 + \beta_1 x_{i1} + \dots + \beta_p x_{ip} + \varepsilon_i,
  \qquad \varepsilon_i \sim \N(0, \sigma^2).
\end{equation}
%
Mô hình giả định quan hệ là tuyến tính (sau khi biến đổi biến nếu cần), các sai
số độc lập, cùng phương sai và có phân phối chuẩn. Biến giả $\USA$ bằng 1 với xe
Mỹ và 0 với xe nước ngoài, nên hệ số của nó là chênh lệch hệ số chặn giữa hai
nhóm. Số hạng tương tác $\HP \times \USA$ cho phép hệ số góc theo HP khác nhau
giữa hai nhóm. Khi hai biến giải thích tương quan mạnh, hệ số của chúng không ổn
định và khó diễn giải riêng; hiện tượng đa cộng tuyến này được đo bằng hệ số
phóng đại phương sai VIF.

Phương pháp bình phương tối thiểu (OLS) cho một giá trị ước lượng và một khoảng
tin cậy theo tần suất, còn cách tiếp cận Bayes cho cả phân phối hậu nghiệm của
mỗi hệ số. Với tiên nghiệm ít thông tin, trung bình hậu nghiệm gần bằng ước
lượng OLS; ví dụ hệ số $\beta_{\HP}$ của M3 là $-0{,}860$ theo OLS và $-0{,}859$
theo Bayes.

\section{Chọn biến}

City là biến phản hồi và HP là biến giải thích chính, đúng với câu hỏi Q4.
Weight và Engine Size tương quan mạnh với HP (\num{0.74} và \num{0.73}) và với
nhau (\num{0.85}), nên nhóm tính VIF trước khi đưa cả ba vào
(bảng~\ref{tab:vif}). Mọi VIF đều dưới 5, nhưng Engine Size gần như lặp lại thông
tin của Weight, nên nhóm chỉ giữ Weight. Với bộ biến của M4, VIF lớn nhất là
\num{2.34}.

\begin{table}[tp]
  \centering
  \caption{Hệ số phóng đại phương sai của hai bộ biến giải thích.}
  \label{tab:vif}
  \begin{tabular}{lrr}
    \toprule
    Biến & HP + Weight + Engine Size & M4: HP + Weight + USA \\
    \midrule
    HP          & \num{2.42} & \num{2.24} \\
    Weight      & \num{3.93} & \num{2.34} \\
    Engine Size & \num{3.84} & -- \\
    USA         & --         & \num{1.07} \\
    \bottomrule
  \end{tabular}
\end{table}

\section{Biến đổi biến}
\label{sec:transform}

Quan hệ giữa City và HP có dạng cong, nên nhóm so sánh hai cách: hồi quy OLS City
theo HP và hồi quy $\log(\City)$ theo HP. Để so được hai mô hình, logarit hàm hợp
lý của mô hình logarit được tính lại trên thang City bằng cách cộng thêm Jacobian
$-\sum_i \log \City_i$. Theo bảng~\ref{tab:transform}, $\log(\City)$ khớp tốt hơn
rõ rệt, AIC thấp hơn 36 đơn vị, và đường làm trơn của phần dư ở
hình~\ref{fig:transform} không còn dạng cong.

\begin{table}[tp]
  \centering
  \caption{So sánh hai thang của biến phản hồi trong hồi quy OLS theo HP;
    logarit hàm hợp lý và AIC tính trên thang City.}
  \label{tab:transform}
  \begin{tabular}{lrrr}
    \toprule
    Biến phản hồi & $R^2$ & Logarit hàm hợp lý & AIC \\
    \midrule
    City         & \num{0.452} & \num{-264.0} & \num{534.0} \\
    $\log(\City)$ & \num{0.506} & \num{-246.0} & \num{498.0} \\
    \bottomrule
  \end{tabular}
\end{table}

\resultfig[0.95]{fig_3_3_response_scale}{Phần dư OLS theo giá trị khớp, với City
  và $\log(\City)$ làm biến phản hồi; đường cam là đường làm trơn
  LOWESS.}{fig:transform}

Vì vậy mọi mô hình đều dùng $y = \log(\City)$ đã chuẩn hoá. HP và Weight được
chuẩn hoá bằng cách trừ trung bình và chia độ lệch chuẩn mẫu, còn USA giữ nguyên
0/1. Chuẩn hoá giúp đặt tiên nghiệm dễ hơn vì mọi hệ số cùng một thang, và giúp
MCMC hội tụ tốt hơn. Kết quả cuối được đổi ngược về MPG bằng
$\City = \exp(\bar{\ell} + s_\ell\, y)$, trong đó $\bar{\ell}$ và $s_\ell$ là
trung bình và độ lệch chuẩn của $\log(\City)$ trong mẫu.

\section{Năm mô hình ứng viên}
\label{sec:candidates}

Năm mô hình ở bảng~\ref{tab:models} thêm dần biến giải thích. M5 dùng cùng bộ
biến với M4, mô hình tốt nhất trong M1--M4 theo PSIS-LOO, nhưng thay hàm hợp lý
chuẩn bằng Student-t để giảm ảnh hưởng của điểm ngoại lai.

\begin{table}[tp]
  \centering
  \caption{Năm mô hình ứng viên; biến phản hồi luôn là $\log(\City)$ đã chuẩn
    hoá.}
  \label{tab:models}
  \begin{tabular}{lll}
    \toprule
    Mô hình & Công thức & Hàm hợp lý \\
    \midrule
    M1 & City $\sim$ HP                        & Chuẩn \\
    M2 & City $\sim$ HP + USA                  & Chuẩn \\
    M3 & City $\sim$ HP + USA + HP $\times$ USA & Chuẩn \\
    M4 & City $\sim$ HP + Weight + USA         & Chuẩn \\
    M5 & City $\sim$ HP + Weight + USA         & Student-t \\
    \bottomrule
  \end{tabular}
\end{table}

\section{Đặc tả đầy đủ cho M3}

Trên thang đã chuẩn hoá, M3 có dạng
%
\begin{equation}
  \label{eq:m3}
  \begin{aligned}
    y_i &\sim \N(\mu_i, \sigma^2), \\
    \mu_i &= \beta_0 + \beta_1 \HP_i + \beta_2 \USA_i + \beta_3 \HP_i \USA_i, \\
    \beta_j &\sim \N(0, 2{,}5^2), \qquad j = 0, \dots, 3, \\
    \sigma &\sim \text{Half-Normal}(1).
  \end{aligned}
\end{equation}
%
Các mô hình khác bỏ bớt hoặc thay số hạng trong $\mu_i$. M5 thay hàm hợp lý chuẩn
bằng $y_i \sim t_\nu(\mu_i, \sigma)$ với $\nu \sim \text{Gamma}(2;\, 0{,}1)$, cho
phép phân phối có đuôi dày hơn khi dữ liệu có điểm ngoại lai. Tiên nghiệm này có
trung bình 20 nhưng vẫn cho xác suất đáng kể với các giá trị $\nu$ nhỏ.

\section{Chọn tiên nghiệm và kiểm tra dự đoán tiên nghiệm}
\label{sec:priors}

Nhóm so sánh ba tiên nghiệm cho các hệ số: tiên nghiệm phẳng $\N(0, 100^2)$, tiên
nghiệm ít thông tin $\N(0, 2{,}5^2)$ và tiên nghiệm hẹp $\N(0, 0{,}5^2)$. Tiên
nghiệm liên hợp Normal--Inverse-Gamma cho công thức hậu nghiệm ở
mục~\ref{sec:conjugate}; còn các tiên nghiệm độc lập trong \eqref{eq:m3} không
liên hợp nên phải dùng MCMC. Để kiểm tra, nhóm sinh \num{1000} bộ dữ liệu từ tiên
nghiệm của M3 và đếm tỉ lệ City nằm trong miền hợp lý. Vì mô hình đặt trên thang
logarit nên City luôn dương, điều kiện còn lại là City không quá 80 MPG.

\begin{table}[tp]
  \centering
  \caption{Kiểm tra dự đoán tiên nghiệm cho M3 với ba tiên nghiệm, mỗi tiên
    nghiệm \num{1000} lần mô phỏng.}
  \label{tab:prior-pred}
  \begin{tabular}{lrrr}
    \toprule
    Tiên nghiệm & City $\le 80$ MPG & City trong 5--80 MPG & Trung vị (MPG) \\
    \midrule
    Phẳng, sd $= 100$        & \num{53.3}\% & \num{3.4}\%  & \num{5.4} \\
    Ít thông tin, sd $= 2{,}5$ & \num{92.1}\% & \num{86.2}\% & \num{21.0} \\
    Hẹp, sd $= 0{,}5$       & \num{99.9}\% & \num{99.9}\% & \num{21.6} \\
    \bottomrule
  \end{tabular}
\end{table}

Theo bảng~\ref{tab:prior-pred} và hình~\ref{fig:prior-pred}, tiên nghiệm phẳng
sinh ra phần lớn giá trị vô lý, chỉ 3{,}4\% mô phỏng nằm trong khoảng 5--80 MPG.
Tiên nghiệm ít thông tin giữ 86\% mô phỏng trong khoảng đó mà vẫn đủ rộng để
không ép kết quả, nên được chọn làm tiên nghiệm chính. Tiên nghiệm hẹp cho giá
trị hợp lý nhưng có thể kéo hệ số về 0; mục~\ref{sec:sensitivity} kiểm tra điều
này.

\resultfig[0.7]{fig_3_6_prior_predictive}{Phân phối City mô phỏng từ tiên nghiệm
  của M3 với ba tiên nghiệm, so với dữ liệu thật; trục hoành theo thang logarit,
  đường đứt đỏ ở 80 MPG.}{fig:prior-pred}

\section{Ảnh hưởng của tiên nghiệm đến kết quả}
\label{sec:sensitivity}

Nhóm chạy lại M3 với cả ba tiên nghiệm. Tiên nghiệm phẳng và tiên nghiệm ít
thông tin cho hậu nghiệm gần như trùng nhau, trung bình lệch nhau nhiều nhất
\num{0.04} độ lệch chuẩn hậu nghiệm, tức kết quả do dữ liệu quyết định. Tiên
nghiệm hẹp kéo các hệ số về 0 nhiều nhất \num{0.42} độ lệch chuẩn hậu nghiệm, rõ
nhất ở tương tác $\HP \times \USA$: hệ số này giảm từ \num{0.294} xuống
\num{0.239} và HDI 95\% bắt đầu chứa 0 (hình~\ref{fig:sensitivity}). Vì vậy kết
luận về tương tác ở mục~\ref{sec:q2} cần nói thận trọng.

\resultfig[0.7]{fig_3_7_sensitivity}{Hậu nghiệm của M3 với ba tiên nghiệm, trung
  bình và HDI 95\%.}{fig:sensitivity}

% ===========================================================================
\chapter{Quá trình ước lượng Bayes}
\label{sec:estimation}

\section{Công thức hậu nghiệm cho M1}
\label{sec:conjugate}

Theo định lý Bayes, hậu nghiệm tỉ lệ với tích của hàm hợp lý và tiên nghiệm:
%
\begin{equation}
  \label{eq:bayes}
  p(\beta, \sigma^2 \mid y) \propto p(y \mid \beta, \sigma^2)\, p(\beta, \sigma^2).
\end{equation}
%
Với tiên nghiệm liên hợp Normal--Inverse-Gamma $\beta \mid \sigma^2 \sim
\N(m_0, \sigma^2 V_0)$ và $\sigma^2 \sim \text{IG}(a_0, b_0)$, hậu nghiệm có cùng
dạng với các tham số
%
\begin{equation}
  \label{eq:nig}
  \begin{aligned}
    V_n &= \bigl(V_0^{-1} + X\trans X\bigr)^{-1}, &
    m_n &= V_n \bigl(V_0^{-1} m_0 + X\trans y\bigr), \\
    a_n &= a_0 + \tfrac{n}{2}, &
    b_n &= b_0 + \tfrac{1}{2}\bigl(y\trans y + m_0\trans V_0^{-1} m_0 - m_n\trans V_n^{-1} m_n\bigr).
  \end{aligned}
\end{equation}
%
Lấy tích phân theo $\sigma^2$ ta được phân phối biên của $\beta$ là phân phối $t$
nhiều chiều với $2a_n$ bậc tự do, tâm $m_n$ và ma trận tỉ lệ $(b_n / a_n) V_n$
\cite{gelman2013bda}. Nhóm chọn $m_0 = 0$, $V_0 = 2{,}5^2 I$, $a_0 = 2$,
$b_0 = 1$ và tính được $a_n = 48{,}5$, $b_n = 23{,}76$.

Để kiểm tra code, nhóm viết lại đúng tiên nghiệm này trong PyMC và lấy mẫu bằng
NUTS. Bảng~\ref{tab:conjugate} cho thấy hai kết quả khớp đến chữ số thập phân thứ
ba; nếu lệch nhau thì code có lỗi. Dòng thứ ba dùng tiên nghiệm ít thông tin độc
lập của M1, khác tiên nghiệm liên hợp, nhưng hậu nghiệm vẫn gần như trùng, vì cả
hai tiên nghiệm đều rộng và 93 quan sát đủ để dữ liệu quyết định kết quả
(hình~\ref{fig:conjugate}).

\begin{table}[tp]
  \centering
  \caption{Hậu nghiệm của $\beta_{\HP}$ trong M1, thang chuẩn hoá: công thức so
    với MCMC.}
  \label{tab:conjugate}
  \begin{tabular}{lrrr}
    \toprule
    Cách tính & Trung bình & Độ lệch chuẩn & Khoảng 95\% \\
    \midrule
    Công thức (NIG)                   & \num{-0.7101} & \num{0.0737} & [\num{-0.855}; \num{-0.565}] \\
    NUTS, cùng tiên nghiệm NIG        & \num{-0.7099} & \num{0.0742} & [\num{-0.857}; \num{-0.564}] \\
    NUTS, tiên nghiệm ít thông tin M1 & \num{-0.7103} & \num{0.0754} & [\num{-0.862}; \num{-0.560}] \\
    \bottomrule
  \end{tabular}
\end{table}

\resultfig[0.95]{fig_4_1_conjugate}{Hậu nghiệm của M1: mật độ $t$ tính bằng công
  thức (đen), NUTS với cùng tiên nghiệm NIG (xanh) và NUTS với tiên nghiệm của M1
  (đỏ).}{fig:conjugate}

\section{Lấy mẫu MCMC}
\label{sec:mcmc}

M2--M5 không có công thức hậu nghiệm nên được lấy mẫu bằng NUTS
\cite{hoffman2014nuts} trong PyMC 5.28 \cite{abrilpla2023pymc}, với 4 chuỗi, mỗi
chuỗi \num{2000} vòng sau \num{1000} vòng khởi động, seed cố định 2026. NUTS là
một dạng của Hamiltonian Monte Carlo: nó dùng gradient của logarit hậu nghiệm để
đề xuất các bước đi xa mà vẫn có xác suất chấp nhận cao, và tự chọn độ dài bước
thay vì để người dùng chỉnh tay.

Để thấy sự khác biệt, nhóm tự viết thuật toán Metropolis--Hastings bước ngẫu
nhiên cho M1. Thuật toán lấy mẫu trên $(\beta_0, \beta_1, \log\sigma)$, cộng thêm
Jacobian $\log\sigma$ vào logarit hậu nghiệm, bước đề xuất có phân phối chuẩn với
độ lệch chuẩn \num{0.08}, chạy 4 chuỗi \num{6000} vòng và bỏ \num{1000} vòng đầu.
Tỉ lệ chấp nhận từ \num{0.41} đến \num{0.43}. Metropolis--Hastings cho cùng hậu
nghiệm với NUTS ($\beta_{\HP} = -0{,}708$ so với $-0{,}710$) nhưng kém hiệu quả
hơn nhiều: cỡ mẫu hiệu dụng khoảng \num{1700} từ \num{20000} mẫu, trong khi NUTS
đạt khoảng \num{12000} từ \num{8000} mẫu. Lý do là mỗi bước của
Metropolis--Hastings chỉ dịch đi một đoạn ngắn, nên các mẫu liên tiếp tương quan
mạnh (hình~\ref{fig:mh}).

\resultfig[0.85]{fig_4_2_mh_trace}{Biểu đồ vết của Metropolis--Hastings tự viết
  cho M1, bốn chuỗi.}{fig:mh}

\section{Chẩn đoán hội tụ}

Nhóm dùng bốn tiêu chí: $\Rhat$ (tính theo thứ hạng) \cite{vehtari2021rhat} nhỏ
hơn \num{1.01}, cỡ mẫu hiệu dụng (ESS) bulk và tail đều trên 400, không có bước
phân kỳ (divergence), và biểu đồ vết cùng biểu đồ thứ hạng không có dấu hiệu bất
thường. Cả năm mô hình đều đạt (bảng~\ref{tab:diagnostics}); biểu đồ vết của mô
hình cuối ở hình~\ref{fig:trace} cho thấy bốn chuỗi trộn đều. Code sẽ dừng nếu có
mô hình không đạt, nên không có hệ số nào được lấy từ chuỗi chưa hội tụ.

\begin{table}[tp]
  \centering
  \caption{Chẩn đoán hội tụ của năm mô hình; giá trị xấu nhất trên mọi tham
    số.}
  \label{tab:diagnostics}
  \begin{tabular}{lrrrr}
    \toprule
    Mô hình & $\Rhat$ lớn nhất & ESS bulk nhỏ nhất & ESS tail nhỏ nhất & Số bước phân kỳ \\
    \midrule
    M1 & \num{1.00} & \num{10830} & \num{5460} & 0 \\
    M2 & \num{1.00} & \num{6013}  & \num{5295} & 0 \\
    M3 & \num{1.00} & \num{5618}  & \num{5197} & 0 \\
    M4 & \num{1.00} & \num{5116}  & \num{5091} & 0 \\
    M5 & \num{1.00} & \num{4626}  & \num{4763} & 0 \\
    \bottomrule
  \end{tabular}
\end{table}

\resultfig[0.85]{fig_4_3_trace_M5}{Biểu đồ vết của M5: mật độ hậu nghiệm theo
  chuỗi (trái) và giá trị qua các vòng lặp (phải).}{fig:trace}

\section{Bảng tóm tắt hậu nghiệm}
\label{sec:posterior}

Bảng~\ref{tab:posterior} cho trung bình hậu nghiệm và HDI 95\% trên thang chuẩn
hoá. Hệ số của HP và Weight là mức thay đổi của $\log(\City)$ chuẩn hoá khi biến
tăng một độ lệch chuẩn; hệ số của USA là chênh lệch giữa xe Mỹ và xe nước ngoài.
M5 có thêm $\nu = 9{,}4$ với HDI 95\% là $[1{,}8;\, 24{,}7]$, tức đuôi phân phối
dày vừa phải.

\begin{table}[tp]
  \centering
  \small
  \caption{Trung bình hậu nghiệm và HDI 95\% của các hệ số, thang chuẩn hoá. Cột
    $\sigma$ là độ lệch chuẩn sai số với M1--M4 và tham số thang đo Student-t với
    M5.}
  \label{tab:posterior}
  \begin{tabular}{lccccc}
    \toprule
    Mô hình & $\beta_{\HP}$ & $\beta_{\USA}$ & $\beta_{\HP \times \USA}$ & $\beta_{\Wt}$ & $\sigma$ \\
    \midrule
    M1 & \num{-0.71} [\num{-0.86}; \num{-0.56}] & -- & -- & -- & \num{0.72} \\
    M2 & \num{-0.70} [\num{-0.84}; \num{-0.55}] & \num{-0.41} [\num{-0.69}; \num{-0.13}] & -- & -- & \num{0.69} \\
    M3 & \num{-0.86} [\num{-1.07}; \num{-0.65}] & \num{-0.40} [\num{-0.69}; \num{-0.14}] & \num{0.29} [\num{0.03}; \num{0.59}] & -- & \num{0.68} \\
    M4 & \num{-0.15} [\num{-0.29}; \num{-0.00}] & \num{-0.16} [\num{-0.36}; \num{0.03}] & -- & \num{-0.76} [\num{-0.89}; \num{-0.61}] & \num{0.47} \\
    M5 & \num{-0.11} [\num{-0.25}; \num{0.02}] & \num{-0.15} [\num{-0.33}; \num{0.02}] & -- & \num{-0.76} [\num{-0.89}; \num{-0.62}] & \num{0.38} \\
    \bottomrule
  \end{tabular}
\end{table}

Khi thêm Weight ở M4 và M5, hệ số của HP giảm từ khoảng $-0{,}7$ xuống khoảng
$-0{,}1$, còn $\sigma$ giảm từ \num{0.68} (M3) xuống \num{0.47} (M4). Xe mạnh
thường cũng là xe nặng, nên liên hệ giữa HP và City trong M1--M3 phần lớn chính là
liên hệ giữa Weight và City. Dữ liệu quan sát không đủ để tách riêng tác động của
từng biến (hình~\ref{fig:forest}). Lưu ý ở M1--M4, $\sigma$ là độ lệch chuẩn của
sai số; còn ở M5, $\sigma$ là tham số thang đo của phân phối Student-t nên không
so trực tiếp được với các mô hình khác (xem mục~\ref{sec:q3}).

\resultfig[0.7]{fig_4_4_forest}{Hệ số hậu nghiệm của M1--M5 trên thang chuẩn
  hoá, trung bình và HDI 95\%.}{fig:forest}

% ===========================================================================
\chapter{Đánh giá kết quả và trả lời câu hỏi}
\label{sec:evaluation}

\section{Kiểm tra dự đoán hậu nghiệm}

Với mỗi mô hình, nhóm sinh \num{8000} bộ dữ liệu từ hậu nghiệm, đổi về MPG, rồi
so mật độ của chúng với dữ liệu thật (hình~\ref{fig:ppc}). Cả năm mô hình đều tái
tạo tốt trung bình của City, với p-value Bayes từ \num{0.38} đến \num{0.47}.
M1--M4 không tái tạo được hai đầu phân phối: dữ liệu sinh ra hiếm khi có xe đạt
46 MPG như Geo Metro, p-value của giá trị lớn nhất chỉ từ \num{0.006} đến
\num{0.024}. M5 tốt hơn ở cả hai đầu, p-value của giá trị lớn nhất là \num{0.041}
và của giá trị nhỏ nhất là \num{0.044}.

\resultfig[0.95]{fig_5_1_ppc}{Kiểm tra dự đoán hậu nghiệm: mật độ dữ liệu thật
  (đen) và 100 bộ dữ liệu sinh từ hậu nghiệm (xanh).}{fig:ppc}

\section{So sánh mô hình}
\label{sec:compare}

Nhóm so sánh các mô hình bằng ELPD ước lượng qua PSIS-LOO và WAIC
\cite{vehtari2017loo}, tính bằng ArviZ \cite{kumar2019arviz}. Hai tiêu chí cho
cùng thứ hạng (bảng~\ref{tab:compare}). Mô hình được chọn là M5, có ELPD cao nhất
và trọng số stacking \num{0.97}. M5 hơn M4 \num{4.45} đơn vị ELPD, với sai số
chuẩn \num{2.60}, tức khoảng \num{1.7} lần sai số chuẩn: đủ để chọn M5 nhưng
chênh lệch không lớn. Nhóm có Weight (M4, M5) hơn nhóm không có (M1--M3) từ 34
đến 43 đơn vị ELPD, tức hơn bốn lần sai số chuẩn.

\begin{table}[tp]
  \centering
  \caption{So sánh mô hình bằng PSIS-LOO và WAIC; chênh lệch tính so với M5.}
  \label{tab:compare}
  \begin{tabular}{lrrrrrr}
    \toprule
    Mô hình & ELPD-LOO & $p_{\text{loo}}$ & Chênh lệch & SE chênh lệch & ELPD-WAIC & Pareto $k > 0{,}5$ \\
    \midrule
    M5 & \num{-59.8}  & \num{6.4} & 0           & --         & \num{-59.7}  & 0 \\
    M4 & \num{-64.2}  & \num{6.7} & \num{4.4}   & \num{2.6}  & \num{-64.2}  & 1 \\
    M3 & \num{-98.2}  & \num{6.0} & \num{38.4}  & \num{8.4}  & \num{-98.1}  & 0 \\
    M2 & \num{-99.3}  & \num{5.1} & \num{39.6}  & \num{8.2}  & \num{-99.3}  & 0 \\
    M1 & \num{-102.4} & \num{4.1} & \num{42.6}  & \num{7.8}  & \num{-102.4} & 0 \\
    \bottomrule
  \end{tabular}
\end{table}

Không có quan sát nào có Pareto $k$ vượt \num{0.7}, nên ước lượng PSIS-LOO dùng
được cho mọi mô hình (hình~\ref{fig:compare}).

\resultfig[0.7]{fig_5_2_compare}{ELPD-LOO của năm mô hình, kèm sai số chuẩn và
  sai số chuẩn của chênh lệch so với mô hình tốt nhất.}{fig:compare}

\section{Trả lời Q1: quan hệ giữa các biến}
\label{sec:q1}

Công suất và hiệu suất nhiên liệu có quan hệ nghịch rõ rệt, nhưng khi thêm trọng
lượng vào mô hình thì hệ số của công suất nhỏ đi nhiều. Trên thang logarit, khi
HP tăng 10 thì trung vị có điều kiện của City được nhân với
$\exp(10\, \beta_{\HP}\, s_\ell / s_{\HP})$, trong đó $s_{\HP}$ là độ lệch chuẩn
của HP. Vì vậy các phần trăm dưới đây là mức thay đổi tương đối của trung vị
City.

\begin{itemize}
  \item Chỉ dùng HP (M1): mỗi 10 hp tăng thêm làm City giảm 3{,}0\%, HDI 95\% từ
    $-3{,}6\%$ đến $-2{,}4\%$. Từ 140 lên 150 hp, trung vị City giảm \num{0.66}
    MPG, HDI 95\% từ $-0{,}80$ đến $-0{,}52$ MPG (hình~\ref{fig:q1}).
  \item Với cùng Weight (M5): mức giảm ứng với 10 hp chỉ còn $-0{,}45\%$, HDI 95\%
    từ $-1{,}07\%$ đến $+0{,}10\%$, xác suất hệ số âm là \num{0.94}.
  \item Với cùng HP và xuất xứ: mỗi \SI{100}{kg} nặng hơn ứng với trung vị City
    thấp hơn 6{,}2\%, HDI 95\% từ $-7{,}1\%$ đến $-5{,}0\%$.
\end{itemize}

Các biểu đồ phân tán ở mục~\ref{sec:eda} cũng cho thấy điều này: City giảm theo
cả HP, Weight và Engine Size theo dạng cong, và tương quan với Weight mạnh hơn
với HP. City và Open Road có tương quan \num{0.94}, nên Open Road nhiều khả năng
có xu hướng tương tự; tuy nhiên muốn có kết luận định lượng cho Open Road thì
phải ước lượng mô hình riêng cho biến này.

\resultfig[0.7]{fig_5_3_q1_m1_fit}{M1: City theo HP, với HDI 95\% của trung vị
  có điều kiện (xanh đậm) và khoảng dự đoán 95\% (xanh nhạt).}{fig:q1}

\section{Trả lời Q2: xe Mỹ và xe nước ngoài}
\label{sec:q2}

Xe Mỹ kém tiết kiệm hơn xe nước ngoài cùng công suất, nhưng khoảng cách này giảm
nhiều khi thêm trọng lượng vào mô hình, vì xe Mỹ trong mẫu nặng hơn
(bảng~\ref{tab:q2}). Điều này cho thấy xuất xứ và trọng lượng gắn chặt với nhau
trong dữ liệu, chứ chưa chứng minh trọng lượng là nguyên nhân. Khi chỉ xét cùng
HP, xe Mỹ có City thấp hơn khoảng 8{,}6\% và xác suất hậu nghiệm để xe Mỹ tiết
kiệm hơn chỉ là \num{0.003}. Khi xét thêm cùng Weight, khoảng cách chỉ còn
3{,}3\% và HDI 95\% chạm 0.

\begin{table}[tp]
  \centering
  \caption{Hệ số của USA trong ba mô hình; chênh lệch City đổi ra phần trăm.}
  \label{tab:q2}
  \begin{tabular}{lrrr}
    \toprule
    Mô hình & $\beta_{\USA}$ [HDI 95\%] & Chênh lệch City của xe Mỹ & $P(\beta_{\USA} > 0 \mid y)$ \\
    \midrule
    M2: HP + USA         & \num{-0.41} [\num{-0.69}; \num{-0.13}] & $-8{,}7\%$ & \num{0.003} \\
    M3: có HP $\times$ USA & \num{-0.40} [\num{-0.69}; \num{-0.14}] & $-8{,}6\%$ & \num{0.003} \\
    M5: có Weight        & \num{-0.15} [\num{-0.33}; \num{0.02}]  & $-3{,}3\%$ & \num{0.043} \\
    \bottomrule
  \end{tabular}
\end{table}

Trong M3, hệ số tương tác $\HP \times \USA$ dương, bằng \num{0.29} với HDI 95\%
$[0{,}03;\, 0{,}59]$ và xác suất dương \num{0.98}. Nghĩa là xe nước ngoài tiết
kiệm hơn hẳn ở công suất thấp nhưng mất dần lợi thế khi HP tăng, và hai đường hồi
quy gặp nhau quanh 200 hp (hình~\ref{fig:q2}, trái). Theo mục~\ref{sec:sensitivity},
kết quả này thay đổi khi đổi tiên nghiệm, nên nhóm chỉ coi đây là gợi ý chứ chưa
phải kết luận chắc chắn. Khi giữ Weight ở mức trung bình (hình~\ref{fig:q2},
phải), hai đường gần như trùng nhau.

\resultfig[0.95]{fig_5_4_q2_usa}{Đường hồi quy và HDI 95\% cho xe Mỹ và xe nước
  ngoài: M3 với tương tác (trái) và M5 với Weight cố định ở mức trung bình
  (phải).}{fig:q2}

\section{Trả lời Q3: các điểm ngoại lai}
\label{sec:q3}

Nhóm xét phần dư chuẩn hoá và Pareto $k$ từ M4 (bản hàm hợp lý chuẩn của mô hình
cuối), rồi so với M5 (bảng~\ref{tab:outliers}, hình~\ref{fig:q3}). Phần dư chuẩn
hoá là hiệu giữa $y_i$ và trung bình hậu nghiệm của $\mu_i$, chia cho trung bình
hậu nghiệm của $\sigma$. Pareto $k$ cho biết một quan sát ảnh hưởng đến hậu
nghiệm nhiều hay ít: $k$ lớn nghĩa là nếu bỏ quan sát đó thì hậu nghiệm thay đổi
nhiều.

\begin{table}[tp]
  \centering
  \small
  \caption{Các xe được xét là ngoại lai, theo M4.}
  \label{tab:outliers}
  \begin{tabular}{lrrrr>{\raggedright\arraybackslash}p{6.4cm}}
    \toprule
    Xe & City & Dự đoán & Phần dư & $k$ & Giải thích \\
    \midrule
    Honda Civic        & 42 & \num{28.0} & \num{3.86}  & \num{0.54}
      & Xe nhẹ (\SI{1066}{kg}), có thể là phiên bản động cơ tiết kiệm nhiên liệu; ảnh hưởng lớn nhất \\
    Geo Metro          & 46 & \num{34.9} & \num{2.65}  & \num{0.36}
      & Xe nhẹ nhất mẫu (\SI{769}{kg}), máy \SI{1.0}{L} 55 hp, nằm ngoài vùng có nhiều dữ liệu \\
    Suzuki Swift       & 39 & \num{32.0} & \num{1.90}  & \num{0.21}
      & Cùng loại xe cỡ nhỏ với Geo Metro, máy \SI{1.3}{L}, \SI{891}{kg} \\
    Mazda RX-7         & 17 & \num{21.8} & \num{-2.37} & \num{0.29}
      & Động cơ rotary \SI{1.3}{L} nhưng 255 hp; thân nhẹ nên mô hình dự đoán cao hơn \\
    Volkswagen Corrado & 18 & \num{23.4} & \num{-2.52} & \num{0.19}
      & Xe thể thao, máy \SI{2.8}{L} 178 hp trên thân \SI{1275}{kg} \\
    Ford Aerostar      & 15 & \num{17.7} & \num{-1.57} & \num{0.10}
      & Xe van nặng; bất thường khi chỉ dùng HP, hết bất thường khi có Weight \\
    Mercedes-Benz 300E & 19 & \num{18.6} & \num{0.19}  & \num{0.09}
      & Ngoại lai về giá (Max Price 80, cao nhất), không phải về nhiên liệu \\
    \bottomrule
  \end{tabular}
\end{table}

Nhóm xe van có City từ 15 đến 18 MPG trông bất thường khi mô hình chỉ dùng HP,
nhưng khi có Weight thì phần dư của chúng chỉ từ $-1{,}6$ đến $+0{,}3$. Vậy City
thấp của các xe này là do chúng nặng chứ không phải do công suất. Ba xe cỡ nhỏ
Honda Civic, Geo Metro và Suzuki Swift tiết kiệm hơn mức mà một đường
log-tuyến tính dự đoán, còn Mazda RX-7 và Volkswagen Corrado tốn nhiên liệu hơn
dự đoán. Giải thích về phiên bản động cơ của Honda Civic chỉ là giả thuyết của
nhóm, vì dữ liệu không ghi phiên bản.

So với M4, M5 giảm trọng số của các điểm này: Pareto $k$ lớn nhất giảm từ
\num{0.54} xuống \num{0.46}, và $\beta_{\HP}$ giảm độ lớn từ $-0{,}15$ xuống
$-0{,}11$, HDI 95\% giờ chứa 0. Hệ số của Weight giữ nguyên ở $-0{,}76$, nên liên
hệ giữa trọng lượng và City không phụ thuộc vào các điểm ngoại lai.

Tham số $\sigma$ của M5 bằng \num{0.38}, nhỏ hơn \num{0.47} của M4, nhưng hai số
này không so trực tiếp được. Với Student-t, $\sigma$ là tham số thang đo; độ lệch
chuẩn của sai số bằng $\sigma\sqrt{\nu/(\nu-2)}$ khi $\nu > 2$ và vô hạn khi
$\nu \le 2$ (chỉ \num{0.3}\% mẫu hậu nghiệm có $\nu \le 2$). Khi đổi về cùng
thước đo, độ lệch chuẩn sai số của M5 có trung vị hậu nghiệm khoảng \num{0.45},
gần bằng M4. Vậy M5 không làm sai số nhỏ đi; khác biệt là phần lớn xe nằm sát
đường hồi quy hơn, còn vài xe ngoại lai rơi vào phần đuôi dày của Student-t.

\resultfig[0.95]{fig_5_5_q3_outliers}{Phần dư chuẩn hoá theo giá trị dự đoán của
  M4 (trái) và Pareto $k$ của từng quan sát với M4 và M5 (phải).}{fig:q3}

\section{Trả lời Q4: dự đoán cho xe 300 hp}
\label{sec:q4}

Biến phản hồi là City và biến giải thích chính là Horsepower. Nếu chỉ biết công
suất thì M1 là mô hình đúng với câu hỏi của đề: trung vị City của một xe 300 hp
khoảng \num{13.6} MPG, khoảng dự đoán 95\% từ \num{9.7} đến \num{18.8} MPG. M5 cần
biết thêm xuất xứ và trọng lượng, nên kết quả của M5 chỉ áp dụng cho một xe cụ
thể. Với xe Mỹ 300 hp nặng \SI{1681}{kg} (trọng lượng trung bình của ba xe thật
có từ 295 hp trở lên), trung vị City khoảng \num{16.6} MPG, khoảng dự đoán 95\%
từ \num{13.4} đến \num{20.6} MPG (bảng~\ref{tab:q4}). Một xe 300 hp nặng hơn hay
nhẹ hơn sẽ có dự đoán khác.

Vì mô hình đặt trên $\log(\City)$, $\exp(\mu)$ là trung vị có điều kiện của City
chứ không phải trung bình; cột thứ ba của bảng là trung bình hậu nghiệm của trung
vị này. Với M5, trung bình của City khi $\log(\City)$ có phân phối Student-t
thậm chí là vô hạn, nên nhóm chỉ báo cáo trung vị và khoảng dự đoán. Khoảng dự
đoán khác HDI của trung vị ở chỗ nó tính cả sai số của từng xe, nên rộng hơn
nhiều.

\begin{table}[tp]
  \centering
  \caption{Dự đoán City cho một xe 300 hp, \SI{1681}{kg} nếu mô hình có Weight.}
  \label{tab:q4}
  \begin{tabular}{llrr}
    \toprule
    Mô hình & Xuất xứ & Trung vị City [HDI 95\%] & Khoảng dự đoán 95\% \\
    \midrule
    M1: chỉ HP & bất kỳ      & \num{13.6} [\num{12.2}; \num{15.0}] & \num{9.7}--\num{18.8} \\
    M3         & Mỹ          & \num{14.3} [\num{12.5}; \num{16.1}] & \num{10.3}--\num{19.6} \\
    M3         & nước ngoài  & \num{12.9} [\num{10.9}; \num{14.8}] & \num{9.2}--\num{17.9} \\
    M5         & Mỹ          & \num{16.6} [\num{15.4}; \num{18.0}] & \num{13.4}--\num{20.6} \\
    M5         & nước ngoài  & \num{17.2} [\num{15.9}; \num{18.5}] & \num{13.8}--\num{21.3} \\
    \bottomrule
  \end{tabular}
\end{table}

Ba xe thật gần 300 hp trong mẫu đều là xe Mỹ. Khi dùng thông số riêng của từng
xe, khoảng dự đoán của M5 đều chứa giá trị thật (bảng~\ref{tab:q4-real}). Tuy
nhiên ba xe này nằm trong dữ liệu dùng để ước lượng mô hình, nên đây chỉ là kiểm
tra trên mẫu học, chưa phải kiểm tra khả năng dự đoán cho xe mới. M1 dự đoán thấp
hơn thực tế (\num{13.6} MPG), vì đường log-tuyến tính theo HP tiếp tục đi xuống ở
phía phải, trong khi xe 300 hp thực tế không nặng hơn xe 200 hp bao nhiêu
(hình~\ref{fig:q4}). Vì 300 hp là giá trị lớn nhất trong mẫu, mọi dự đoán ở đây
đều nằm ở rìa dữ liệu.

\begin{table}[tp]
  \centering
  \caption{Dự đoán của M5 cho ba xe thật có từ 295 hp trở lên.}
  \label{tab:q4-real}
  \begin{tabular}{lrrrrr}
    \toprule
    Xe & HP & Weight (kg) & City thật & Trung vị dự đoán & Khoảng dự đoán 95\% \\
    \midrule
    Cadillac Seville   & 295 & \num{1785} & 16 & \num{15.6} & \num{12.6}--\num{19.2} \\
    Chevrolet Corvette & 300 & \num{1533} & 17 & \num{18.3} & \num{14.6}--\num{22.7} \\
    Dodge Stealth      & 300 & \num{1726} & 18 & \num{16.1} & \num{13.0}--\num{20.0} \\
    \bottomrule
  \end{tabular}
\end{table}

\resultfig[0.75]{fig_5_6_q4_300hp}{Phân phối dự đoán hậu nghiệm cho một xe Mỹ
  300 hp theo M1 và M5; đường đỏ là City của ba xe thật gần 300 hp.}{fig:q4}

\section{Hạn chế}

Mẫu chỉ có 93 xe và chỉ có ba xe trên 280 hp, nên dự đoán ở rìa dữ liệu kém chắc
chắn. Dữ liệu từ năm 1993, công nghệ động cơ khi đó khác xa hiện nay, nên kết quả
không áp dụng cho xe ngày nay. Các quan hệ tìm được là tương quan chứ không phải
nhân quả, và dự đoán ngoài khoảng 55--300 hp hoặc 769--\SI{1862}{kg} không đáng
tin. Phần kiểm tra dự đoán ở Q4 chỉ dùng các xe đã có trong dữ liệu học; nhóm
chưa đánh giá khả năng dự đoán trên xe ngoài mẫu.

% ===========================================================================
\chapter{Kết luận}
\label{sec:conclusion}

Trong mẫu 93 xe năm 1993, trọng lượng là biến liên hệ mạnh nhất với hiệu suất
nhiên liệu trong thành phố. Đây là kết luận về liên hệ trên dữ liệu quan sát,
không phải về quan hệ nhân quả. Mô hình cuối M5, $\log(\City) \sim \HP + \Wt +
\USA$ với hàm hợp lý Student-t, hội tụ tốt, tái tạo dữ liệu tốt nhất và đứng đầu
theo cả PSIS-LOO lẫn WAIC.

\paragraph{Q1, quan hệ giữa các biến.} Nếu chỉ xét HP, mỗi 10 hp cao hơn ứng với
trung vị City thấp hơn 3{,}0\%. Với cùng Weight, mức này gần như bằng 0, chỉ còn
$-0{,}45\%$, trong khi mỗi \SI{100}{kg} nặng hơn ứng với trung vị City thấp hơn
6{,}2\%. Open Road tương quan \num{0.94} với City nên nhiều khả năng có xu hướng
tương tự, nhưng nhóm chưa xây dựng mô hình riêng cho biến này.

\paragraph{Q2, xe Mỹ và xe nước ngoài.} Cùng công suất, xe Mỹ kém tiết kiệm hơn
8{,}6\%, xác suất điều ngược lại chỉ là \num{0.003}. Khi tính cả Weight, khoảng
cách chỉ còn 3{,}3\% và HDI 95\% chạm 0. Vậy phần lớn chênh lệch giữa hai nhóm đi
kèm với chênh lệch trọng lượng, dù dữ liệu quan sát không cho phép khẳng định
trọng lượng là nguyên nhân.

\paragraph{Q3, điểm ngoại lai.} Honda Civic, Geo Metro và Suzuki Swift tiết kiệm
bất thường, còn Mazda RX-7 (động cơ rotary) và Volkswagen Corrado tốn nhiên liệu
hơn dự đoán. Xe van chỉ là ngoại lai khi bỏ qua trọng lượng. Các kết luận chính
không đổi khi chuyển sang hàm hợp lý Student-t.

\paragraph{Q4, xe 300 hp.} Nếu chỉ biết công suất (M1), trung vị City của một xe
300 hp khoảng \num{13.6} MPG, khoảng dự đoán 95\% từ \num{9.7} đến \num{18.8} MPG.
Nếu biết thêm đó là xe Mỹ nặng \SI{1681}{kg} (M5), trung vị City khoảng
\num{16.6} MPG, khoảng dự đoán 95\% từ \num{13.4} đến \num{20.6} MPG. Ba xe thật
gần 300 hp đều nằm trong khoảng của M5, nhưng chúng thuộc dữ liệu học nên đây chỉ
là kiểm tra trên mẫu học.

\paragraph{Hướng phát triển.} Có thể dùng mô hình phân cấp theo Type, trong đó
mỗi loại xe có hệ số chặn riêng và các loại xe dùng chung thông tin với nhau, để
tách phần chênh lệch do loại xe. Cũng có thể thêm các biến như Engine Size, Drive
train hoặc Manual, hoặc thử biến phản hồi gallon trên dặm $1/\City$, vì trong
phân tích sơ bộ biến này còn khớp tốt hơn.

% ===========================================================================
\appendix

\chapter{Code và dữ liệu}
\label{sec:appendix}

Toàn bộ phân tích nằm trong notebook chung \texttt{bayes\_cars\_1993.ipynb},
ghép từ năm module trong thư mục \texttt{analysis/}. Notebook chạy từ đầu đến cuối
trong khoảng 8 phút với Python 3.12, PyMC 5.28 và ArviZ 0.23, seed 2026. Phân
công như ở bảng~\ref{tab:modules}.

\begin{table}[tp]
  \centering
  \caption{Các module trong notebook chung.}
  \label{tab:modules}
  \begin{tabular}{lll}
    \toprule
    Module & Phụ trách & Mục trong báo cáo \\
    \midrule
    \texttt{02\_data\_eda.py}   & TV2 & Chương~\ref{sec:data}, mục~\ref{sec:transform} \\
    \texttt{03\_models.py}      & TV3 & Mục~\ref{sec:candidates}--\ref{sec:sensitivity} \\
    \texttt{04\_sampling.py}    & TV4 & Mục~\ref{sec:mcmc}--\ref{sec:posterior} \\
    \texttt{01\_conjugate.py}   & TV1 & Mục~\ref{sec:conjugate} \\
    \texttt{05\_evaluation.py}  & TV5 & Chương~\ref{sec:evaluation} \\
    \bottomrule
  \end{tabular}
\end{table}

Dữ liệu đã làm sạch nằm ở \texttt{outputs/data\_clean.csv}, bảng hậu nghiệm đầy
đủ của mọi mô hình ở \texttt{outputs/tables/tab\_4\_4\_posterior.csv}, và các
hình ở \texttt{outputs/figures/}. Dưới đây là hàm khai báo chung cho năm mô hình
trong PyMC.

\begin{lstlisting}
def build_model(d, sc, terms, prior="weak", likelihood="normal"):
    X = design(d, sc, terms)          # standardised HP, Weight; USA as 0/1
    y = sc.y(d["city"])               # standardised log(City)
    coords = {"term": terms, "obs": d["name"].to_numpy()}
    with pm.Model(coords=coords) as model:
        Xd = pm.Data("X", X, dims=("obs", "term"))
        b0 = pm.Normal("b0", 0, PRIOR_SD[prior])
        b = pm.Normal("b", 0, PRIOR_SD[prior], dims="term")
        sigma = pm.HalfNormal("sigma", 1)
        mu = pm.Deterministic("mu", b0 + pm.math.dot(Xd, b), dims="obs")
        if likelihood == "student":
            nu = pm.Gamma("nu", 2, 0.1)
            pm.StudentT("y", nu=nu, mu=mu, sigma=sigma, observed=y, dims="obs")
        else:
            pm.Normal("y", mu=mu, sigma=sigma, observed=y, dims="obs")
    return model
\end{lstlisting}

\chapter{Khai báo sử dụng AI}
\label{sec:ai}

Nhóm có dùng trợ lý AI \textbf{Claude} của Anthropic trong quá trình làm bài, cụ
thể:

\begin{itemize}
  \item hỗ trợ viết, chú thích và sửa lỗi code Python (PyMC, ArviZ) trong thư mục
    \texttt{analysis/};
  \item hỗ trợ soạn, sửa văn bản và trình bày LaTeX của báo cáo;
  \item rà soát cách diễn giải thống kê, ví dụ phân biệt trung vị với trung bình
    trên thang log, ý nghĩa tham số thang đo của Student-t và tránh diễn giải
    nhân quả.
\end{itemize}

Mọi con số trong báo cáo đều lấy từ kết quả chạy code của nhóm
(\texttt{outputs/}), không lấy từ câu trả lời của AI. Nhóm đã đối chiếu các nhận
định với kết quả tính toán và tài liệu tham khảo, và chịu trách nhiệm về toàn bộ
nội dung bài.

\newpage

\printbibliography[heading=bibintoc,title={Tài liệu tham khảo}]

\end{document}
